% GENERATED FILE. Edit canonical lessons, then run npm run generate:books.
% canonical-source-hash: fnv1a64:1d6bd2f63c506003
% canonical-lessons: BN-C537-monosthir-kora, BN-C537-ogotya, BN-C537-nimraji, BN-C537-onggikar, BN-C537-pichhpa-hooya

\chapter{Making Up My Mind}
\label{ch:bn-making-up-my-mind}
\hlchaptermodality{537}

\begin{chapteropening}
\textbf{By the end of this chapter:} \emph{I can say how I made up my mind, or had no choice.}

Complete the last lesson of chapter 537: I can say how I made up my mind, or had no choice.
\end{chapteropening}

\section[\texorpdfstring{\bn{মনস্থির} \bn{করা} (mônôsthir kôrā)}{মনস্থির করা (mônôsthir kôrā)}]{\bn{মনস্থির} \bn{করা} (mônôsthir kôrā) — to make up one's mind}

\label{lesson:BN-C537-monosthir-kora}



\begin{quote}
Before the new one: say the Bengali for a style, then the Bengali for to fit.
\end{quote}

\subsection*{You'll want to know: \bn{মনস্থির} \bn{করা}}
\textbf{\bn{মনস্থির} \bn{করা}} — \emph{mônôsthir kôrā} — \textquotedblleft{}to make up one's mind\textquotedblright{}.

\begin{quote}
\bn{অনেক} \bn{ভেবে} \bn{আমি} \bn{মনস্থির} \bn{করলাম}: \bn{নতুন} \bn{চাকরিটা} \bn{নেব।}
\end{quote}

\begin{quote}
\emph{Ônek bhebe āmi mônôsthir kôrlām: nôtun chākriṭā nebô.}
\end{quote}

\begin{quote}
After much thought I made up my mind: I'll take the new job.
\end{quote}

Past (he/she …): \bn{তিনি} \bn{শেষে} \bn{মনস্থির} \bn{করলেন}, \emph{tini sheshe mônôsthir kôrlen}.

\begin{cousinweb}
\textbf{\bn{মনস্থির} \bn{করা}} (\emph{mônôsthir kôrā}) joins \bn{মনস্থির}, from Sanskrit \emph{manas}, \textquotedblleft{}mind\textquotedblright{}, and \emph{sthira}, \textquotedblleft{}firm\textquotedblright{} (Monier-Williams), to \bn{করা}, \textquotedblleft{}to do\textquotedblright{}: to make the mind firm.
\end{cousinweb}

\begin{grammarlens}[title={What you've built}]
One more word for saying what happens if.
\end{grammarlens}

\subsection*{Your turn}
\begin{itemize}
\raggedright
  \item \emph{Say it:} \emph{mônôsthir kôrā}
  \item \emph{Say it:} \emph{mônôsthir kôrā}, once more
  \item \emph{Say it:} \emph{Ônek bhebe āmi mônôsthir kôrlām: nôtun chākriṭā nebô.} — After much thought I made up my mind: I'll take the new job.
  \item \emph{Recall:} say the Bengali for a style, then the Bengali for to fit, then say \emph{mônôsthir kôrā} again
\end{itemize}

\begin{tcolorbox}[breakable,colback=teal!4,colframe=teal!35!black,title={Before you move on}]
What does \bn{মনস্থির} \bn{করা} mean? (\textquotedblleft{}To make up one's mind\textquotedblright{}.) Say it once more.
\end{tcolorbox}
\section[\texorpdfstring{\bn{অগত্যা} (ôgôtyā)}{অগত্যা (ôgôtyā)}]{\bn{অগত্যা} (ôgôtyā) — with no other choice; perforce}

\label{lesson:BN-C537-ogotya}



\begin{quote}
Before the new one: say the Bengali for to fit, then the Bengali for to make up one's mind.
\end{quote}

\subsection*{You'll want to know: \bn{অগত্যা}}
\textbf{\bn{অগত্যা}} — \emph{ôgôtyā} — \textquotedblleft{}with no other choice; perforce\textquotedblright{}.

\begin{quote}
\bn{বাস} \bn{এল} \bn{না}, \bn{অগত্যা} \bn{আমরা} \bn{হেঁটে} \bn{গেলাম।}
\end{quote}

\begin{quote}
\emph{Bās elô nā, ôgôtyā āmrā hẽṭe gelām.}
\end{quote}

\begin{quote}
The bus didn't come, so we had no choice but to walk.
\end{quote}

\begin{cousinweb}
\textbf{\bn{অগত্যা}} (\emph{ôgôtyā}) is Sanskrit \emph{agatyā}, \textquotedblleft{}unavoidably\textquotedblright{} (Monier-Williams), literally \textquotedblleft{}for lack of a way\textquotedblright{}.
\end{cousinweb}

\begin{grammarlens}[title={What you've built}]
One more word for saying what happens if.
\end{grammarlens}

\subsection*{Your turn}
\begin{itemize}
\raggedright
  \item \emph{Say it:} \emph{ôgôtyā}
  \item \emph{Say it:} \emph{ôgôtyā}, once more
  \item \emph{Say it:} \emph{Bās elô nā, ôgôtyā āmrā hẽṭe gelām.} — The bus didn't come, so we had no choice but to walk.
  \item \emph{Recall:} say the Bengali for to fit, then the Bengali for to make up one's mind, then say \emph{ôgôtyā} again
\end{itemize}

\begin{tcolorbox}[breakable,colback=teal!4,colframe=teal!35!black,title={Before you move on}]
What does \bn{অগত্যা} mean? (\textquotedblleft{}With no other choice; perforce\textquotedblright{}.) Say it once more.
\end{tcolorbox}
\section[\texorpdfstring{\bn{নিমরাজি} (nimrāji)}{নিমরাজি (nimrāji)}]{\bn{নিমরাজি} (nimrāji) — half-willing; agreeing reluctantly}

\label{lesson:BN-C537-nimraji}



\begin{quote}
Before the new one: say the Bengali for to make up one's mind, then the Bengali for with no other choice.
\end{quote}

\subsection*{You'll want to know: \bn{নিমরাজি}}
\textbf{\bn{নিমরাজি}} — \emph{nimrāji} — \textquotedblleft{}half-willing; agreeing reluctantly\textquotedblright{}.

\begin{quote}
\bn{বাবা} \bn{প্রথমে} \bn{নিমরাজি} \bn{ছিলেন}, \bn{পরে} \bn{রাজি} \bn{হলেন।}
\end{quote}

\begin{quote}
\emph{Bābā prôthôme nimrāji chhilen, pôre rāji hôlen.}
\end{quote}

\begin{quote}
At first my father was half-willing; later he agreed.
\end{quote}

\begin{cousinweb}
\textbf{\bn{নিমরাজি}} (\emph{nimrāji}) joins \bn{নিম}, \textquotedblleft{}half\textquotedblright{}, usually traced to Persian \emph{nīm}, to \bn{রাজি}, \textquotedblleft{}willing\textquotedblright{}: only half willing.
\end{cousinweb}

\begin{grammarlens}[title={What you've built}]
One more word for saying what happens if.
\end{grammarlens}

\subsection*{Your turn}
\begin{itemize}
\raggedright
  \item \emph{Say it:} \emph{nimrāji}
  \item \emph{Say it:} \emph{nimrāji}, once more
  \item \emph{Say it:} \emph{Bābā prôthôme nimrāji chhilen, pôre rāji hôlen.} — At first my father was half-willing; later he agreed.
  \item \emph{Recall:} say the Bengali for to make up one's mind, then the Bengali for with no other choice, then say \emph{nimrāji} again
\end{itemize}

\begin{tcolorbox}[breakable,colback=teal!4,colframe=teal!35!black,title={Before you move on}]
What does \bn{নিমরাজি} mean? (\textquotedblleft{}Half-willing; agreeing reluctantly\textquotedblright{}.) Say it once more.
\end{tcolorbox}
\section[\texorpdfstring{\bn{অঙ্গীকার} (ônggīkār)}{অঙ্গীকার (ônggīkār)}]{\bn{অঙ্গীকার} (ônggīkār) — a pledge; a commitment}

\label{lesson:BN-C537-onggikar}



\begin{quote}
Before the new one: say the Bengali for with no other choice, then the Bengali for half-willing.
\end{quote}

\subsection*{You'll want to know: \bn{অঙ্গীকার}}
\textbf{\bn{অঙ্গীকার}} — \emph{ônggīkār} — \textquotedblleft{}a pledge; a commitment\textquotedblright{}.

\begin{quote}
\bn{নেতা} \bn{গ্রামে} \bn{নতুন} \bn{স্কুলের} \bn{অঙ্গীকার} \bn{করলেন।}
\end{quote}

\begin{quote}
\emph{Netā grāme nôtun skuler ônggīkār kôrlen.}
\end{quote}

\begin{quote}
The leader pledged a new school for the village.
\end{quote}

\begin{cousinweb}
\textbf{\bn{অঙ্গীকার}} (\emph{ônggīkār}) is a tatsama, Sanskrit \emph{aṅgīkāra}, \textquotedblleft{}agreement, a promise\textquotedblright{} (Monier-Williams), literally \textquotedblleft{}making a thing one's own\textquotedblright{}.
\end{cousinweb}

\begin{grammarlens}[title={What you've built}]
One more word for saying what happens if.
\end{grammarlens}

\subsection*{Your turn}
\begin{itemize}
\raggedright
  \item \emph{Say it:} \emph{ônggīkār}
  \item \emph{Say it:} \emph{ônggīkār}, once more
  \item \emph{Say it:} \emph{Netā grāme nôtun skuler ônggīkār kôrlen.} — The leader pledged a new school for the village.
  \item \emph{Recall:} say the Bengali for with no other choice, then the Bengali for half-willing, then say \emph{ônggīkār} again
\end{itemize}

\begin{tcolorbox}[breakable,colback=teal!4,colframe=teal!35!black,title={Before you move on}]
What does \bn{অঙ্গীকার} mean? (\textquotedblleft{}A pledge; a commitment\textquotedblright{}.) Say it once more.
\end{tcolorbox}
\section[\texorpdfstring{\bn{পিছপা} \bn{হওয়া} (pichhpā hôoyā)}{পিছপা হওয়া (pichhpā hôoyā)}]{\bn{পিছপা} \bn{হওয়া} (pichhpā hôoyā) — to shrink back; to hold back}

\label{lesson:BN-C537-pichhpa-hooya}



\begin{quote}
Before the new one: say the Bengali for half-willing, then the Bengali for a commitment.
\end{quote}

\subsection*{You'll want to know: \bn{পিছপা} \bn{হওয়া}}
\textbf{\bn{পিছপা} \bn{হওয়া}} — \emph{pichhpā hôoyā} — \textquotedblleft{}to shrink back; to hold back\textquotedblright{}.

\begin{quote}
\bn{বিপদে} \bn{ছেলেটা} \bn{পিছপা} \bn{হল} \bn{না।}
\end{quote}

\begin{quote}
\emph{Bipôde chheleṭā pichhpā hôlô nā.}
\end{quote}

\begin{quote}
The boy didn't shrink back in the face of danger.
\end{quote}

Past (he/she …): \bn{তিনি} \bn{পিছপা} \bn{হলেন} \bn{না}, \emph{tini pichhpā hôlen nā}.

\begin{cousinweb}
\textbf{\bn{পিছপা} \bn{হওয়া}} (\emph{pichhpā hôoyā}) joins \bn{পিছপা}, \textquotedblleft{}foot backwards\textquotedblright{}, to \bn{হওয়া}, \textquotedblleft{}to become\textquotedblright{}. The first part belongs with \emph{paśca}, \textquotedblleft{}the hind part\textquotedblright{}, under which CDIAL lists Bengali \emph{pichā}, \textquotedblleft{}the rear\textquotedblright{} (CDIAL 7990); the second is \bn{পা}, \textquotedblleft{}a foot\textquotedblright{}.
\end{cousinweb}

\begin{grammarlens}[title={What you've built}]
One more word for saying what happens if.
\end{grammarlens}

\subsection*{Your turn}
\begin{itemize}
\raggedright
  \item \emph{Say it:} \emph{pichhpā hôoyā}
  \item \emph{Say it:} \emph{pichhpā hôoyā}, once more
  \item \emph{Say it:} \emph{Bipôde chheleṭā pichhpā hôlô nā.} — The boy didn't shrink back in the face of danger.
  \item \emph{Recall:} say the Bengali for half-willing, then the Bengali for a commitment, then say \emph{pichhpā hôoyā} again
\end{itemize}

\begin{tcolorbox}[breakable,colback=teal!4,colframe=teal!35!black,title={Before you move on}]
What does \bn{পিছপা} \bn{হওয়া} mean? (\textquotedblleft{}To shrink back; to hold back\textquotedblright{}.) Say it once more.
\end{tcolorbox}
